\chapter{THỰC NGHIỆM VÀ ĐÁNH GIÁ KẾT QUẢ}
\label{chap:thucnghiem}

\section{Cấu hình thực nghiệm}
\label{sec:exp_config}

Toàn bộ thực nghiệm dùng mẫu hm500k sau lọc $k$-core ($k=10$): 7.519 người dùng, 10.345 sản phẩm, chia theo mốc thời gian chung (mục~\ref{subsec:split_sampling}). Tinh chỉnh và đánh giá cuối được chạy trên GPU NVIDIA GeForce RTX 3050 Laptop (4 GB), RAM 16 GB, bằng PyTorch \parencite{paszke2019pytorch}; độ trễ suy diễn được đo riêng trên CPU (mục~\ref{subsec:latency}). Mọi bảng số liệu, hình và con số kết quả trong chương này được sinh tự động từ tệp kết quả của quy trình thực nghiệm, không nhập tay.

Giao thức đánh giá:
\begin{itemize}
    \item \textbf{Full Ranking (chính):} với mỗi người dùng kiểm thử, xếp hạng toàn bộ 10.216 sản phẩm có trong tập huấn luyện $\cup$ xác thực, trừ các sản phẩm người dùng đã mua trong hai tập này; các sản phẩm đúng luôn nằm trong tập ứng viên. Hòa điểm được phá bằng một khoá tất định (seed 2026), giống nhau cho mọi mô hình.
    \item \textbf{Độ đo chính:} NDCG@10 trung bình theo người dùng trên tập kiểm thử (2.996 người dùng, 9.229 cặp đúng). Độ đo phụ: Recall@10, HR@10, Precision@10, NDCG@5 (định nghĩa cho trường hợp nhiều sản phẩm đúng ở mục~\ref{sec:metrics}).
    \item \textbf{Sampled-99 (chỉ đối chiếu):} mỗi cặp đúng được xếp cùng 99 sản phẩm âm lấy ngẫu nhiên, theo giao thức NCF gốc \parencite{he2017ncf}; không dùng để kết luận \parencite{krichene2020sampled}.
    \item \textbf{Số lần chạy:} mỗi mô hình được chạy với 3 seed (42, 2024, 2025); với mỗi seed, số epoch được chọn trên tập xác thực rồi mô hình được huấn luyện lại trên tập huấn luyện $\cup$ xác thực trước khi chấm test (mục~\ref{subsec:split_sampling}); báo cáo trung bình $\pm$ độ lệch chuẩn qua seed; kiểm định theo mục~\ref{sec:protocol}.
    \item \textbf{Lần chấm:} số liệu dưới đây là lần chấm test thứ hai, sau khi bổ sung bước huấn luyện lại; lần chấm đầu cho cùng kết luận (mục~\ref{sec:protocol}, phần lệch kế hoạch).
\end{itemize}

\section{Phương pháp đối chứng (Baselines)}
\label{sec:baselines}

Bảy mô hình được đánh giá, chia ba nhóm:
\begin{enumerate}
    \item \textbf{Baseline:} \textbf{Random} (điểm ngẫu nhiên -- cận dưới); \textbf{Most Popular} (điểm là số người dùng đã mua sản phẩm trong tập huấn luyện, mọi người dùng nhận cùng một danh sách trừ sản phẩm đã mua); \textbf{BPR-MF} \parencite{rendle2009bpr} -- nhân tử hóa ma trận tuyến tính tối ưu xếp hạng theo cặp, được tinh chỉnh cùng ngân sách như NeuMF.
    \item \textbf{Ablation:} \textbf{GMF} (chỉ nhánh MF) và \textbf{MLP} (chỉ nhánh DNN, mô hình DNN thuần -- mục~\ref{subsec:early_fusion_design}).
    \item \textbf{Mô hình đề tài:} \textbf{NeuMF-Scratch} (khởi tạo ngẫu nhiên) và \textbf{NeuMF-Pretrained} (khởi tạo từ GMF, MLP đã huấn luyện) \parencite{he2017ncf} -- mô hình lai MF + DNN theo kiểu Early Fusion ở mức biểu diễn (mục~\ref{subsec:early_late_fusion}).
\end{enumerate}
Lọc cộng tác dựa trên độ tương đồng sản phẩm (ItemKNN \cite{sarwar2001item}) không được đưa vào vì ma trận tương đồng dày $10.345 \times 10.345$ vượt ngưỡng bộ nhớ an toàn của cài đặt; đây là một hạn chế vì Ferrari Dacrema và cộng sự \parencite{dacrema2019progress} cho thấy ItemKNN được tune tốt có thể vượt NeuMF.

\section{Kết quả tinh chỉnh trên tập xác thực}
\label{sec:tuning_results}

Bảng~\ref{tab:tuning} liệt kê cấu hình tốt nhất của từng mô hình sau 6 cấu hình thử (tổng thời gian tinh chỉnh 1,9 giờ GPU, mọi dòng nhật ký cùng một phiên bản mã nguồn, không có thay đổi chưa lưu). Mọi con số trong bảng là trên tập \emph{xác thực}; chúng lạc quan vì chính tập này được dùng để chọn cấu hình, nên chỉ dùng để chọn, không dùng để kết luận.

\bangketqua{tab_tuning.tex}{\texttt{scripts/10\_tune.py}}

Ba nhận xét từ quá trình tinh chỉnh. Thứ nhất, trên validation NeuMF-Scratch đứng đầu (NDCG@10 = \Val{Scr}{NDCG10}), trong khi NeuMF-Pretrained (\Val{Pre}{NDCG10}) chỉ ngang BPR-MF (\Val{BPR}{NDCG10}) -- pre-training không giúp trên dữ liệu này. Theo quy tắc đã đăng ký (mục~\ref{sec:protocol}), \textbf{NeuMF-Scratch là mô hình lai được chọn}. Thứ hai, các mô hình nơ-ron đạt đỉnh rất sớm (đa số ở epoch 1--3, muộn nhất epoch 9), dấu hiệu quá khớp nhanh trên dữ liệu thưa. Thứ ba, các ứng viên đã bị loại khỏi phạm vi trước khi chấm test (mục~\ref{sec:protocol}) đều cao hơn NeuMF-Scratch trên validation: trộn điểm iALS + NeuMF-Scratch đạt 0,01260, DeepCF/CFNet 0,01169 và iALS đã tune 0,01141 (iALS mặc định chỉ 0,00906) -- phù hợp với nhận định của \textcite{rendle2022ials} rằng MF được tune kỹ là baseline rất mạnh.

\section{Kết quả trên tập kiểm thử}
\label{sec:exp_results}

\subsection{So sánh các độ đo xếp hạng (Full Ranking)}
\label{subsec:hr_ndcg_comparison}

Bảng~\ref{tab:final} và Hình~\ref{fig:final_metrics} trình bày kết quả trên tập kiểm thử của 7 mô hình, trung bình qua 3 seed.

\bangketqua{tab_final.tex}{\texttt{scripts/11\_final.py} và \texttt{scripts/12\_significance.py}}

\hinhketqua{media/figures/final/final_metrics.png}{\textwidth}{NDCG@10, Recall@10 và HR@10 trên tập kiểm thử: trung bình $\pm$ độ lệch chuẩn qua 3 seed, chấm đen là từng seed}{fig:final_metrics}{chạy \texttt{scripts/13\_plot\_final.py} rồi \texttt{scripts/15\_export\_report.py}}

Trước hết, thang giá trị tuyệt đối rất thấp (NDCG@10 cao nhất khoảng 0,01): mỗi người dùng kiểm thử có trung bình 3,08 sản phẩm đúng giữa hơn 10.000 ứng viên, và dự đoán những món \emph{chưa từng mua} trong 8 tuần sau mốc huấn luyện là bài toán khó với dữ liệu thời trang. Vì vậy thông tin quan trọng là thứ hạng tương đối và chênh lệch có ý nghĩa thống kê hay không, không phải độ lớn tuyệt đối.

Bốn quan sát chính:
\begin{enumerate}
    \item \textbf{BPR-MF có trung bình cao nhất trên mọi độ đo} (NDCG@10 = \Res{BPR}{NDCG10} $\pm$ \Sd{BPR}{NDCG10}), dù là mô hình tuyến tính đơn giản nhất trong nhóm có học.
    \item \textbf{NeuMF-Pretrained ngang Most Popular}: \Res{Pre}{NDCG10} so với \Res{MostPop}{NDCG10} (chênh \Sig{PreMostPop}{rel}, không có ý nghĩa thống kê). Mô hình lai được chọn theo validation, NeuMF-Scratch, chỉ đạt \Res{Scr}{NDCG10}.
    \item \textbf{GMF, MLP và hai biến thể NeuMF nằm sát nhau} (NDCG@10 từ \Res{GMF}{NDCG10} đến \Res{Pre}{NDCG10}), chênh lệch nhỏ hơn hoặc cùng cỡ với độ lệch chuẩn giữa các seed của các mô hình nơ-ron.
    \item \textbf{Mọi mô hình có học đều vượt xa Random} (NDCG@10 = \Res{Random}{NDCG10}), tức đều học được tín hiệu thật từ dữ liệu. Most Popular là baseline mạnh: độ phổ biến giải thích phần lớn những gì các mô hình phức tạp hơn nắm bắt được.
\end{enumerate}

Ở K = 20 (Bảng~\ref{tab:k20}) -- chỉ số mô tả bổ sung, tính lại từ hạng đã lưu của cùng lần chấm, không chấm lại mô hình -- BPR-MF vẫn có trung bình cao nhất (NDCG@20 = \Kx{BPR}{NDCG20}); các mô hình còn lại tiếp tục nằm sát nhau và thứ tự giữa chúng thay đổi so với K = 10 (Most Popular \Kx{MostPop}{NDCG20} và MLP \Kx{MLP}{NDCG20} nhỉnh hơn NeuMF-Pretrained \Kx{Pre}{NDCG20}), thêm một dấu hiệu rằng chênh lệch giữa các mô hình này nằm trong mức nhiễu.

\bangketqua{tab_k20.tex}{\texttt{scripts/16\_extra\_k.py}}

\subsection{Kiểm định ý nghĩa thống kê}
\label{subsec:multiseed_results}

Bảng~\ref{tab:significance} và Hình~\ref{fig:final_significance} trình bày kết quả kiểm định cặp theo người dùng (\Sig{all}{nusers} người dùng) trên họ \Sig{all}{n} so sánh đã đăng ký trước.

\bangketqua{tab_significance.tex}{\texttt{scripts/12\_significance.py}}

\hinhketqua{media/figures/final/final_significance.png}{0.95\textwidth}{Chênh lệch NDCG@10 (A $-$ B) với khoảng tin cậy 95\% bootstrap theo người dùng, kèm giá trị $p$ đã hiệu chỉnh Holm}{fig:final_significance}{chạy \texttt{scripts/13\_plot\_final.py} rồi \texttt{scripts/15\_export\_report.py}}

Số so sánh đạt tiêu chí ``A tốt hơn B'' là \textbf{\Sig{all}{nbetter}/\Sig{all}{n}}: mọi khoảng tin cậy đều chứa 0 và mọi giá trị $p$ sau hiệu chỉnh Holm đều $\ge 0{,}05$. So sánh gần ngưỡng nhất là NeuMF-Pretrained với GMF: NDCG@10 của NeuMF-Pretrained chênh \Sig{PreGMF}{rel} so với GMF và $p$ Wilcoxon chưa hiệu chỉnh là \Sig{PreGMF}{pw} -- dưới 0,05 nếu xét riêng lẻ -- nhưng sau hiệu chỉnh cho 8 so sánh $p_{Holm} = $ \Sig{PreGMF}{pholm} và khoảng tin cậy [\Sig{PreGMF}{cilo}; \Sig{PreGMF}{cihi}] vẫn chứa 0; đây là ví dụ trực tiếp cho lý do phải hiệu chỉnh đa so sánh. So với BPR-MF, NeuMF-Pretrained chênh \Sig{PreBPR}{rel} ($p_{Holm} = $ \Sig{PreBPR}{pholm}). Kết luận đúng theo tiêu chí đã đăng ký: \emph{NeuMF không vượt BPR-MF, không vượt Most Popular, và không khác biệt có ý nghĩa với các nhánh GMF, MLP của chính nó.}

Hình~\ref{fig:val_vs_test} đặt NDCG@10 trên validation (lúc tinh chỉnh) cạnh NDCG@10 trên test. Thứ hạng không được giữ nguyên: NeuMF-Scratch đứng đầu trên validation nhưng trên test chỉ xếp thứ năm trong bảy mô hình, dưới cả BPR-MF, NeuMF-Pretrained, MLP và Most Popular. Đây là minh chứng trực tiếp rằng chênh lệch giữa các mô hình nhỏ hơn mức nhiễu của việc chọn trên một tập validation 2.275 người dùng -- và là lý do đề tài không kết luận dựa trên một lần chạy hay trên số liệu validation.

\hinhketqua{media/figures/final/final_val_vs_test.png}{0.8\textwidth}{NDCG@10 trên validation (seed 42, cấu hình tốt nhất) so với trên test (trung bình 3 seed)}{fig:val_vs_test}{chạy \texttt{scripts/13\_plot\_final.py} rồi \texttt{scripts/15\_export\_report.py}}

\subsection{Phân tích bóc tách thành phần (Ablation Study)}
\label{subsec:ablation}

\begin{itemize}
    \item \textbf{Đóng góp của hai nhánh GMF và MLP:} Nếu hợp nhất GMF và MLP tạo ra cộng hưởng như khung NCF kỳ vọng, NeuMF phải vượt cả hai nhánh đơn lẻ. Số liệu không xác nhận điều này: NDCG@10 của NeuMF-Pretrained chênh \Sig{PreGMF}{rel} so với GMF và \Sig{PreMLP}{rel} so với MLP, của NeuMF-Scratch chênh \Sig{ScrGMF}{rel} so với GMF và \Sig{ScrMLP}{rel} so với MLP -- cả bốn so sánh đều không có ý nghĩa thống kê ($p_{Holm}$ lần lượt \Sig{PreGMF}{pholm}, \Sig{PreMLP}{pholm}, \Sig{ScrGMF}{pholm}, \Sig{ScrMLP}{pholm}). Tăng số tham số lên gấp đôi (NeuMF-Scratch có 2,30 triệu tham số so với 1,14 triệu của GMF) không chuyển thành lợi ích đo được.
    \item \textbf{Early Fusion và Late Fusion:} NeuMF hợp nhất nhánh MF và nhánh DNN ở mức biểu diễn (Early Fusion, mục~\ref{subsec:early_late_fusion}); như ý trên, trên tập kiểm thử nó không khác biệt có ý nghĩa với từng nhánh đứng riêng. Late Fusion thuần MF + DNN với cùng thành phần -- trộn điểm GMF và MLP được huấn luyện riêng -- chưa được thực hiện, nên đề tài chưa kết luận được cách hợp nhất nào phù hợp hơn. Phép trộn điểm duy nhất đã thử là trộn iALS với NeuMF-Scratch (một mô hình MF với mô hình lai), trọng số chọn trên tập xác thực: NDCG@10 = 0,01260 trên tập xác thực, cao hơn từng thành phần đứng riêng (iALS 0,01141; NeuMF-Scratch 0,01107), gợi ý NeuMF mang tín hiệu bổ sung cho MF. Mô hình này đã bị loại khỏi phạm vi trước lần chấm test đầu tiên (mục~\ref{sec:protocol}) nên chưa có kết quả trên tập kiểm thử.
    \item \textbf{Hiệu quả của Pre-training:} NDCG@10 test của NeuMF-Pretrained chênh \Sig{PreScr}{rel} so với NeuMF-Scratch, không có ý nghĩa thống kê ($p_{Holm} = $ \Sig{PreScr}{pholm}); trên validation thì ngược lại (Pretrained thấp hơn Scratch). Không có bằng chứng rằng Pre-training giúp trên dữ liệu này -- khác với mức cải thiện 1--2\% mà bài gốc báo cáo trên MovieLens và Pinterest \parencite{he2017ncf}.
    \item \textbf{Kích thước Embedding và tỉ lệ mẫu âm:} Được tinh chỉnh như mọi siêu tham số khác (Bảng~\ref{tab:hyperparams}). GMF và NeuMF-Scratch chọn $d=64$, MLP chọn $d=32$; MLP và NeuMF-Pretrained chọn 8 mẫu âm, các mô hình còn lại chọn 4 (Bảng~\ref{tab:tuning}). Số tầng của tháp MLP được giữ cố định theo thiết kế NCF ($[2d, d, d/2, d/4]$), chưa được khảo sát riêng.
\end{itemize}

\subsection{Đánh giá phân tầng: sản phẩm phổ biến/ít phổ biến và người dùng ít/nhiều tương tác}
\label{subsec:longtail_results}

Để kiểm tra mô hình có học được sở thích cá nhân hay chỉ nắm bắt độ phổ biến, tập kiểm thử được chia theo hai chiều, định nghĩa chỉ dựa trên dữ liệu mà mô hình được huấn luyện trước khi chấm test (tập huấn luyện $\cup$ xác thực): \textit{(i)} theo sản phẩm -- nhóm \textbf{head} là 10\% sản phẩm được mua nhiều nhất, còn lại là \textbf{tail}; \textit{(ii)} theo người dùng -- nhóm \textbf{cold} là 20\% người dùng có ít tương tác nhất (cold-start \emph{tương đối}, không phải người dùng mới), còn lại là \textbf{warm}. Đây là phân tích mô tả, không kiểm định (mục~\ref{sec:protocol}).

\bangketqua{tab_stratified.tex}{\texttt{scripts/14\_secondary.py}}

Trên nhóm sản phẩm tail -- nơi độ phổ biến không giúp được -- NeuMF-Pretrained đạt NDCG@10 \Strat{Pre}{tail}, BPR-MF \Strat{BPR}{tail} và Most Popular \Strat{MostPop}{tail}; trên nhóm head lần lượt là \Strat{Pre}{head}, \Strat{BPR}{head} và \Strat{MostPop}{head}. Trên nhóm người dùng cold, NeuMF-Pretrained đạt \Strat{Pre}{cold} so với \Strat{Pre}{warm} ở nhóm warm. Gần như toàn bộ độ chính xác đến từ nhóm head: trên nhóm này mọi mô hình đạt NDCG@10 khoảng 0,03, còn trên nhóm tail không mô hình nào vượt quá 0,0003 (GMF, MLP và Most Popular bằng 0; cao nhất là NeuMF-Scratch với \Strat{Scr}{tail}). Nói cách khác, không mô hình nào -- kể cả NeuMF -- gợi ý đúng được sản phẩm ít phổ biến. Khoảng cách cold/warm nhỏ và không theo một chiều: Most Popular đạt \Strat{MostPop}{cold} ở nhóm cold, cao nhất trong mọi mô hình và cao hơn chính nó ở nhóm warm (\Strat{MostPop}{warm}); chỉ BPR-MF và NeuMF-Scratch thấp hơn ở nhóm cold (\Strat{BPR}{cold} so với \Strat{BPR}{warm}; \Strat{Scr}{cold} so với \Strat{Scr}{warm}). Như vậy, với người dùng ít lịch sử, gợi ý theo độ phổ biến trên dữ liệu này không kém các mô hình cá nhân hóa -- nhận xét mô tả, chưa kiểm định, và nhóm cold chỉ có \Strat{n}{cold} người dùng.

Bảng~\ref{tab:beyond} bổ sung các chỉ số ngoài độ chính xác của danh sách top-10 (mục~\ref{subsec:beyond_accuracy}).

\bangketqua{tab_beyond.tex}{\texttt{scripts/14\_secondary.py}}

Most Popular là trường hợp giới hạn: mọi người dùng nhận gần như cùng một danh sách gồm các sản phẩm phổ biến nhất, nên độ bao phủ catalog chỉ \Beyond{MostPop}{coverage} và tỉ lệ sản phẩm head trong gợi ý là \Beyond{MostPop}{HRR}. Các mô hình cá nhân hóa có độ bao phủ lần lượt \Beyond{BPR}{coverage} (BPR-MF), \Beyond{GMF}{coverage} (GMF), \Beyond{MLP}{coverage} (MLP), \Beyond{Scr}{coverage} (NeuMF-Scratch) và \Beyond{Pre}{coverage} (NeuMF-Pretrained). NeuMF-Scratch đa dạng nhất: độ bao phủ \Beyond{Scr}{coverage}, độ mới \Beyond{Scr}{novelty} bit và tỉ lệ sản phẩm head thấp nhất (\Beyond{Scr}{HRR}) -- nhưng độ chính xác không cao hơn (NDCG@10 \Res{Scr}{NDCG10}, xếp thứ năm). Ngược lại, MLP gần như chỉ gợi ý sản phẩm phổ biến (độ bao phủ \Beyond{MLP}{coverage}, HRR \Beyond{MLP}{HRR}), sát với Most Popular; NeuMF-Pretrained và GMF cũng có HRR trên 98\%. BPR-MF vừa có độ chính xác trung bình cao nhất vừa đa dạng thứ hai (độ bao phủ \Beyond{BPR}{coverage}, HRR \Beyond{BPR}{HRR}), nên trên dữ liệu này đa dạng hơn không nhất thiết đi kèm độ chính xác thấp hơn.

\subsection{Đối chiếu với giao thức Sampled-99}
\label{subsec:sampled_comparison}

Bảng~\ref{tab:sampled} chấm lại cùng các mô hình theo giao thức của bài NCF gốc: mỗi trong \Samp{all}{ncases} cặp đúng của tập kiểm thử được xếp cùng 99 sản phẩm âm lấy ngẫu nhiên (cố định, dùng chung cho mọi mô hình).

\bangketqua{tab_sampled.tex}{\texttt{scripts/14\_secondary.py}}

Mọi mô hình đều bị thổi phồng chỉ số dưới Sampled-99 (NeuMF-Pretrained: NDCG@10 \Samp{Pre}{NDCG10}, gấp \Samp{Pre}{ratio} lần Full Ranking; Most Popular gấp \Samp{MostPop}{ratio} lần), vì xếp một sản phẩm đúng giữa 100 ứng viên dễ hơn nhiều so với giữa hơn 10.000. Thứ hạng giữa các mô hình \Samp{all}{rankchanged} thay đổi so với Full Ranking: Most Popular chỉ xếp thứ tư theo Full Ranking nhưng đứng đầu theo Sampled-99. Đây là minh họa trực tiếp cho cảnh báo của \textcite{krichene2020sampled}: chỉ số lấy mẫu không nhất quán với chỉ số đầy đủ, nên không được dùng để so sánh mô hình -- và là lý do các con số của bài gốc (HR@10 khoảng 0,7 trên MovieLens) không thể đặt cạnh kết quả của đề tài.

\subsection{Độ trễ suy diễn}
\label{subsec:latency}

Để có bằng chứng định lượng về khả năng triển khai, độ trễ phục vụ một yêu cầu gợi ý được đo trên CPU: với mỗi người dùng chọn ngẫu nhiên trong tập kiểm thử, mô hình chấm toàn bộ khoảng \Lat{all}{ncand} ứng viên rồi lấy top-10 (Bảng~\ref{tab:latency}).

\bangketqua{tab_latency.tex}{\texttt{scripts/14\_secondary.py}}

Độ trễ trung vị của NeuMF-Pretrained là \Lat{Pre}{p50} ms, phân vị 95 là \Lat{Pre}{p95} ms (BPR-MF: \Lat{BPR}{p50} / \Lat{BPR}{p95} ms). Ở quy mô catalog khoảng 10.000 sản phẩm, chấm toàn bộ catalog cho một người dùng trên CPU là khả thi cho phục vụ trực tuyến. Với catalog lớn hơn nhiều bậc, cách chấm toàn bộ bằng MLP sẽ không còn phù hợp: tích vô hướng của MF/BPR cho phép dùng tìm kiếm lân cận gần đúng, còn MLP thì không \parencite{rendle2020ncfmf}.

\section{Phạm vi và giới hạn thực nghiệm}
\label{sec:pham_vi_thuc_nghiem}

Kết quả của chương này có giá trị trong phạm vi sau: \textit{(i)} mẫu hm500k (khoảng 1,6\% khách hàng của H\&M) sau lọc $k=10$, thiên về khách mua nhiều; \textit{(ii)} giao thức chia theo mốc thời gian chung; khi huấn luyện lại trên tập huấn luyện $\cup$ xác thực, số epoch được lấy từ lần huấn luyện chỉ trên tập huấn luyện (ít hơn khoảng 3,5\% dữ liệu) vì không còn tập xác thực để dừng sớm, nên có thể chưa tối ưu; bước huấn luyện lại này được bổ sung sau lần chấm test đầu tiên (mục~\ref{sec:protocol}); \textit{(iii)} chỉ người dùng và sản phẩm đã có trong dữ liệu huấn luyện (tập huấn luyện $\cup$ xác thực khi chấm test) -- bài toán Cold-start tuyệt đối nằm ngoài phạm vi; \textit{(iv)} ngày mua chỉ ở mức ngày, nên một sản phẩm kiểm thử có thể được mua cùng ngày (cùng giỏ hàng) với một sản phẩm trong lịch sử -- ảnh hưởng như nhau với mọi mô hình; \textit{(v)} lọc $k$-core được áp dụng trên toàn bộ mẫu trước khi chia, tức có dùng thông tin số lượng tương tác của giai đoạn xác thực/kiểm thử -- rò rỉ nhỏ và như nhau cho mọi mô hình; \textit{(vi)} một bộ dữ liệu duy nhất. Các mô hình nơ-ron trên GPU khác có thể cho số lệch nhẹ do phép tính cuDNN không tất định; kết luận thống kê mới là điều cần được giữ nguyên khi chạy lại.

\section{Thảo luận}
\label{sec:discussion}

\textit{NeuMF có vượt các baseline hay không?} Không. Trên dữ liệu và giao thức của đề tài, cả hai biến thể NeuMF đều không vượt BPR-MF đã tinh chỉnh (BPR-MF có trung bình cao nhất) và không khác biệt có ý nghĩa với Most Popular. Kết quả này nhất quán với hai nghiên cứu đối chứng đáng tin cậy nhất về NCF: \textcite{rendle2020ncfmf} cho thấy MF tích vô hướng được tune kỹ vượt NeuMF ngay trên dữ liệu của bài gốc, và \textcite{dacrema2019progress} cho thấy nhiều cải thiện của mô hình học sâu biến mất khi baseline được tune đúng và tập test không bị dùng để chọn mô hình. Đề tài áp dụng đúng hai điều kiện đó (cùng ngân sách tinh chỉnh, khoá tập test) và thu được kết luận tương tự.

\textit{Kết hợp tuyến tính -- phi tuyến có mang lại lợi ích?} Không đo được. NeuMF -- hợp nhất hai nhánh ở mức biểu diễn -- không khác biệt có ý nghĩa với GMF hay MLP đứng riêng. Một cách giải thích nhất quán với số liệu: với trung bình 26,9 tương tác huấn luyện mỗi người dùng, dữ liệu không đủ để tầng phi tuyến học được cấu trúc tương tác nào vượt quá những gì tích vô hướng đã nắm bắt -- các mô hình nơ-ron đạt đỉnh trên tập xác thực chỉ sau 1--9 epoch rồi quá khớp. \textcite{rendle2020ncfmf} còn chỉ ra rằng một MLP cần rất nhiều dữ liệu mới xấp xỉ được chính phép tích vô hướng.

\textit{Kết quả âm có giá trị gì?} Kết quả âm được báo cáo đúng như quy tắc đã đăng ký trước, không chọn lọc số liệu có lợi. Giá trị của đề tài nằm ở một quy trình đánh giá đáng tin cậy -- dữ liệu tái lập được, chia theo thời gian không rò rỉ tương lai, cùng ngân sách tinh chỉnh, khoá tập test, kiểm định cặp có hiệu chỉnh đa so sánh -- và ở bằng chứng rằng trên dữ liệu giao dịch thời trang thưa, độ phức tạp của NeuMF không được đền bù bằng độ chính xác. Các hướng có tín hiệu tích cực trên validation (iALS được tune, trộn điểm MF với NeuMF) được chuyển thành hướng phát triển ở Chương~\ref{chap:ketluan}.
